\documentclass[11pt]{article}
\usepackage[margin=1in]{geometry}
\usepackage{amsmath,amssymb,amsthm,hyperref}
\setlength{\emergencystretch}{1em}
\newtheorem{theorem}{Theorem}
\newtheorem{lemma}[theorem]{Lemma}
\newtheorem{remark}[theorem]{Remark}
\newcommand{\T}{\mathbb T}
\newcommand{\Z}{\mathbb Z}
\newcommand{\R}{\mathbb R}
\newcommand{\E}{\mathbb E}
\newcommand{\poly}{\operatorname{poly}}
\newcommand{\Lie}{\operatorname{Lie}}
\newcommand{\Lip}{\operatorname{Lip}}
\title{One-variable factorization and approximate additive linearization}
\author{Author proposal: operator-local agent /root/astra\_v14\_baseline\_compiler}
\date{Protocol 1.4; exact source contracts for the Sections 9--11 repair}
\begin{document}
\maketitle
\noindent This is a new connecting proof, pending independent review. It does not
approve itself, the complete inverse theorem, or any release. The earlier
Section 13 packet is unchanged. We use the final published sources [GTZ],
[GT12], [GT14], and [GT08] identified in the accompanying source record.
Only one-variable factorization is asserted here.

\section{Objects, source differences, and the required outputs}
All nilpotent Lie groups below are standard, connected and simply connected;
all lattices and filtrations are standard and rational. Subgroups described as
rational are connected, and their intersection with the ambient lattice is
cocompact in that subgroup. We use the natural ultrapowers of these objects,
with the internal-function convention of [GTZ, Appendix A]. A bounded quantity
is bounded by some standard constant. In the finite combinatorial proof all
constants depend only on the stated positive density and error constants.
Passing that finite proof to an ultrapower therefore preserves uniformity and
internal choices. Braces denote the representative in $(-1/2,1/2]$.

An $\mathbb N$-filtration $G_\bullet$ may have $G_1\ne G_0=G$.
Its positive groups satisfy $[G_i,G_j]\subseteq G_{i+j}$ and are rational.
The polynomial convention is the ordinary degree filtration on the source.
In fixed adapted Mal'cev coordinates, write the coordinates of a polynomial
map as a finite polynomial in $n$ or $\binom n j$. Coordinate degree is bounded
in terms of the fixed filtered group. For a smooth factor, ``small Taylor
coefficients'' means that its constant coefficient is bounded and its
coefficient of index $j\geq1$ is $O(N^{-j})$, in the group Taylor convention
of [GTZ, Lemma B.9]. Equivalent coordinate bounds are justified below.

The statement of [GTZ, D.2] cannot be used as a literal equivalence: on
$\R/\Z$, $g(n)=n/N$ is equidistributed on $[N]$ but is not equidistributed
on $[\lfloor N/2\rfloor]$. For example its first Fourier average on that
half interval tends to $2i/\pi$, whereas its full-interval average tends to
zero. Its nontrivial horizontal character is $1/N$-Lipschitz. We instead use
the one-variable weak factorization [GT12, Proposition 9.2], whose proof uses
the quantitative theorem on an actual dense progression. The published
correction [GT14, pp.1175--1176] explicitly retains the one-variable results.
Its minor corrections on pp.1182--1183 are part of this source choice.

The properness requirement printed in the second bullet of [GTZ, D.4] is
removed. Only a weak-factorization step that starts with failed total
equidistribution forces a proper subgroup. Also, a rational periodic factor
is not asserted to take values in $\Gamma$.

\begin{theorem}[One-variable factorization, DF-D4]\label{thm:D}
Let $N$ be unbounded and let $g\in{}^*\poly(\Z,G_\bullet)$.
There are a standard connected rational subgroup $P\leq G$ and internal
polynomial maps $\epsilon,p,\gamma\in{}^*\poly(\Z,G_\bullet)$ such that
\[
             g(n)=\epsilon(n)p(n)\gamma(n)\qquad(n\in{}^*\Z).
\]
The map $p$ takes values in ${}^*P$, with its induced filtration. On $[N]$,
$\epsilon$ is bounded, is $O(1/N)$-Lipschitz, and has the small Taylor
coefficient bounds above. There is a standard discrete subgroup
$\Lambda\geq\Gamma$ with $\gamma(n)\in{}^*\Lambda$ for every $n$, and a
standard integer $q>0$ with
$\gamma(n+q)\Gamma=\gamma(n)\Gamma$ for every $n$.
For every standard finite-index subgroup $\Delta\leq P\cap\Gamma$,
$p(n)\Delta$ is totally equidistributed in $P/\Delta$ on $[N]$.
The subgroup $P$ is allowed to equal $G$ or to be trivial.
\end{theorem}

\begin{theorem}[Approximate four-function linearization, DF-F1]\label{thm:F}
Let $N$ be unbounded, let $j\geq1$ be standard, and put $e=N^{-j}$.
Let $H\subseteq[-N,N]\cap{}^*\Z$ be internal, and let
$\xi_i:H\to{}^*\T$, $1\leq i\leq4$, be internal functions.
Suppose that for at least $cN^3$ ordered quadruples in $H^4$ satisfying
$h_1+h_2=h_3+h_4$ one has
\[
 \|\xi_1(h_1)+\xi_2(h_2)+\xi_3(h_3)+\xi_4(h_4)-a\|_\T\leq C e,
\]
where $c>0$ and $C<\infty$ are standard and $a\in{}^*\T$.
For any specified index $i$ there are standard $K<\infty$, $c'>0$,
an internal $H'\subseteq H$ with $|H'|\geq c'N$, internal
$\alpha_k\in{}^*\T$, $\beta_k\in{}^*\R$, and $\delta\in{}^*\T$ such that
\[
 \xi_i(h)=\delta+\sum_{k=1}^K\{\alpha_kh\}\beta_k+O(e)
                    \quad(h\in H').                                      \tag{F}
\]
The implicit constant is standard and uniform over $H'$.
The bracket sum restricted to $H'$ is an exact Freiman homomorphism:
it preserves every relation $h_1+h_2=h_3+h_4$ with all four points in $H'$.
No initial Freiman property, boundedness of the real $\beta_k$, or regularity
of the input functions is assumed. All choices may be made internally.
\end{theorem}

\section{Weak factorization and coefficient bounds}
\begin{lemma}[Weak step, DF-D3]\label{lem:weak}
Suppose first that $G_1=G$. If $g(n)\Gamma$ is not totally equidistributed
on $[N]$, it admits the factorization of Theorem~\ref{thm:D} with $p$ taking
values in a standard proper connected rational subgroup, without the final
equidistribution requirement. All other bounds and the discrete-overgroup
and period assertions hold.
\end{lemma}
\begin{proof}
Failed total equidistribution supplies an internal dense progression and a
standard Lipschitz test with noninfinitesimal discrepancy. Thus for some
standard $0<\delta<1/2$ it fails total $\delta$-equidistribution in the
quantitative convention of [GT12]. Reduce $\delta$ further to include the
rationality constant of the fixed Mal'cev basis. Apply the one-variable case
of [GT12, Proposition 9.2, pp.522--523] with this $\delta$, the actual degree
of $G_\bullet$, and the actual dimension of $G$. Its bounds are standard.
The same argument works for an internal Lipschitz test with a standard norm
bound: a finite uniform net in the standard Lipschitz ball first supplies a
standard test with positive discrepancy. This uses compactness of $G/\Gamma$.

Here are the stronger bounds extracted from that source's construction,
including the distinctions needed for later use. Write $g(0)=c\lambda$
with $c$ in a fixed bounded fundamental set and $\lambda\in{}^*\Gamma$.
Then $\widetilde g=c^{-1}g\lambda^{-1}$ has value $1$ at zero. Left
translation by bounded $c$ changes a Lipschitz norm by a standard factor,
and right multiplication by $\lambda^{-1}$ changes no point of $G/\Gamma$.
Thus the same failed-total-equidistribution argument applies to
$\widetilde g$, with a possibly smaller standard $\delta$.

The source's horizontal-character step gives a nonzero standard integer
vector $k$ with
\[
 \psi(\widetilde g(n))=\sum_{r\geq1}t_r\binom n r,
 \qquad \|k\cdot t_r\|_\T=O(N^{-r}).                         \tag{D1}
\]
Here $k\cdot\psi$ is the real lift of the horizontal homomorphism; in
particular it kills all commutators. Standardness of $k$ follows because the
quantitative bound leaves finitely many integer vectors. The source obtains
(D1) by applying its Theorem 2.9 to the witnessing progression and undoing
that progression by its Lemma 8.4. This use does not assert the false
single-scale equivalence described above.

Choose integers $a_r$ nearest to $k\cdot t_r$.
The coordinate support condition in [GT12, Lemma 6.7] says that $t_r$ lies
in the coordinate subspace for $G_r$. If $k$ vanishes on that subspace,
take $a_r=0$ and $v_r=w_r=0$. Otherwise fix a rational vector $u_r$ in
that subspace with $k\cdot u_r=1$, and put
\[
 v_r=a_ru_r,\qquad w_r=(k\cdot t_r-a_r)u_r.
\]
Define $\epsilon_0$ and $\gamma_0$ by Mal'cev coordinates
\[
 \psi(\epsilon_0(n))=\sum_r w_r\binom n r,
 \qquad \psi(\gamma_0(n))=\sum_r v_r\binom n r,
 \qquad p=\epsilon_0^{-1}\widetilde g\gamma_0^{-1}.             \tag{D2}
\]
The same coordinate-support lemma and closure of polynomial maps under
products show that these are polynomial for the original filtration.
The real horizontal lift vanishes identically on $p$, since its coefficient
at every index is $k\cdot t_r-(k\cdot t_r-a_r)-a_r=0$.
Consequently $p$ lies in the connected kernel of that real homomorphism.
This is a standard proper rational subgroup. A kernel of the circle-valued
homomorphism alone would not establish connectedness; (D2) supplies the real
zero required here.

Each $w_r$ is $O(N^{-r})$. The vectors $v_r$ have a common standard
denominator, but their numerators need not be bounded. This denominator
statement is all that is required. By [GT12, A.11--A.12], $\gamma_0$ is
uniformly rational and periodic modulo $\Gamma$ with a standard period.
The more explicit common discrete overgroup argument below also applies.
Set $\epsilon=c\epsilon_0$ and $\gamma=\gamma_0\lambda$.
These are the correct ordered factors of the original $g$.
Multiplication by bounded $c$ preserves the stated small-coordinate bounds,
and the added $\lambda$ lies in the original lattice. The next lemma verifies
all closure and Taylor assertions used here and in the descent.
\end{proof}

\begin{lemma}[Closure of the controlled factors]\label{lem:closure}
In a fixed rational nilpotent group, finite products of bounded polynomial
maps with coefficient bounds $O(N^{-r})$ retain those bounds and are
$O(1/N)$-Lipschitz on $[N]$. These bounds imply the group Taylor bounds in
Theorem~\ref{thm:D}. A finite collection of rational lattices commensurable
with $\Gamma$, together with sets of points whose rational coordinates have
a fixed common denominator, is contained in a standard discrete subgroup
$\Lambda\geq\Gamma$. Any fixed-degree polynomial map into this group is
periodic modulo $\Gamma$ with a standard period.
\end{lemma}
\begin{proof}
Coordinate products, inverses, and changes between first-kind and Mal'cev
coordinates are fixed polynomials by [GT12, A.2--A.3]. In a coordinate
polynomial $\sum a_r n^r$, the bounds $a_r=O(N^{-r})$ are equivalent to
bounded coefficients after substituting $n=Nt$. Products of such polynomials
retain that property. Converting back to the binomial basis gives at index
$r$ a fixed linear combination of coefficients of monomial indices $\ell\geq r$,
hence $O(N^{-r})$. Constant coefficients stay bounded. On the fixed compact
coordinate range for $0\leq t\leq1$, the group-coordinate maps have bounded
derivatives; differentiation in $n$ contributes $N^{-1}$. This proves the
Lipschitz assertion, with any fixed standard metric.

For clarity about the group Taylor convention, one may assign formal weight
$r$ to each coefficient of $n^r$ in the logarithm of a coordinate polynomial.
The recursive Taylor extraction by $g(0)$, then $g_1^{\binom n1}$, then
$g_2^{\binom n2}$, and so on uses only fixed BCH polynomials. Its coefficient
$\log g_r$ contains no monomial of formal weight less than $r$.
Indeed after removing indices below $r$, each surviving polynomial term of
formal weight $w$ has ordinary degree at most $w$. The residual group map
is the identity at $0,\ldots,r-1$, as a formal identity in its input
coefficients. Its logarithmic terms of weight $w<r$ are therefore polynomials
of degree less than $r$ with $r$ zeros, and vanish. BCH products add weights,
so this assertion is preserved at each extraction step.
The remaining finitely many monomials have weight at least $r$, and are
$O(N^{-r})$. This proves the group Taylor assertion, including under products
and fixed rational embeddings of subgroups.

For the discrete-overgroup assertion, use a rational basis adapted to the
lower central series and let its weights be $w_i\geq1$. All points of each
given lattice have logarithmic coordinates with bounded denominators, as do
the specified coordinate sets, by the fixed polynomial coordinate changes.
Choose a standard integer $D$ divisible by these denominators and by the
denominators of all coefficients in the finite BCH formula in this basis.
Increase $D$ for the finitely many basis conversions if necessary. The set
\[
 L=\sum_i D^{-(2w_i-1)}\Z X_i
\]
is closed under BCH. For a nonlinear term with $t\geq2$ inputs of weights
$w_{i_1},\ldots,w_{i_t}$ contributing to output weight $w_k$, one has
$w_k\geq\sum w_{i_a}$. Therefore the difference of denominator exponents is
\[
 (2w_k-1)-\sum_a(2w_{i_a}-1)\geq t-1\geq1.
\]
The extra power of $D$ clears that BCH coefficient. Linear terms and negation
also preserve $L$. Thus $\exp L$ is a standard discrete subgroup containing
all the required points and $\Gamma$. It is cocompact, and
$[\exp L:\Gamma]$ is finite. This argument also proves commensurability
claims needed when a lattice of a rational subgroup is enlarged inside $G$.

For the period in the required right-coset convention, write the Mal'cev
coordinate vector of $\gamma(n)$ as a polynomial of fixed degree $D_0$.
Its values have a common standard coordinate denominator, since
$\gamma(n)\in{}^*\Lambda$. Interpolating at $0,\ldots,D_0$ shows that
its coefficients have a common standard denominator too. Their numerators
may be arbitrary internal integers. By the rational polynomial formulas
for multiplication and inversion, the vector
\[
             R(n,r)=\psi\bigl(\gamma(n)^{-1}\gamma(n+r)\bigr)
\]
is a polynomial in the two integer variables $n,r$ with a common standard
coefficient denominator $A$. The identity $R(n,0)=0$ is a polynomial
identity, so every monomial has a positive power of $r$.
For the standard integer $q=A$, every coordinate of $R(n,q)$ is an integer
for every internal integer $n$. Adapted Mal'cev coordinates identify
$\Gamma$ with the integer-coordinate points. Hence
$\gamma(n)^{-1}\gamma(n+q)\in{}^*\Gamma$, exactly the assertion
$\gamma(n+q)\Gamma=\gamma(n)\Gamma$.
This proves the conclusion of [GT12, A.12] directly; its printed left
quotient $\gamma(n+q)\gamma(n)^{-1}$ is not the quotient used here.
\end{proof}

\section{Minimal subgroup descent, with all finite covers retained}
\begin{proof}[Proof of Theorem~\ref{thm:D}]
First address the filtration convention. If $G_1$ is proper, write
$g(0)=c\lambda$ as above and set $p=c^{-1}g\lambda^{-1}$.
Then $p(0)=1$ and
\[
 p(n)=\lambda\bigl(g(0)^{-1}g(n)\bigr)\lambda^{-1}\in G_1.
\]
The groups $G_r$ are normal in $G_0$, so conjugation preserves them.
The induced filtration on $G_1$ has zeroth and first group both $G_1$,
and the positive groups remain $G_r$. Hence the one-variable result with
the convention $G_0=G_1$ applies there. This also shows directly that a
controlled representation exists: $\epsilon=c$, $\gamma=\lambda$,
and the represented subgroup is $G_1$. If $G_1=G$, the identity-factor
representation is already enough. Trivial groups cause no exception.

Call a subgroup $P$ a representing subgroup if a factorization with all
conclusions of Theorem~\ref{thm:D} except total equidistribution exists.
Choose one with least dimension among the nonempty collection of standard
rational representing subgroups. This uses a least element of the finite
set $\{0,\ldots,\dim G\}$, not an infinite iteration or a uniform choice
over all finite covers.

Suppose that for a standard finite-index $\Delta\leq P\cap\Gamma$ the
orbit $p\Delta$ is not totally equidistributed. The induced filtration of
$P$ is rational for $\Delta$, since passing to a finite-index lattice
does not change its rational Lie algebra. If $P\cap G_1$ is proper in $P$,
the normalization just given already represents $p$ in that proper subgroup,
with a bounded constant on the left and a $\Delta$-element on the right.
Otherwise apply Lemma~\ref{lem:weak} in $P/\Delta$ to obtain
\[
                 p=\epsilon_2p_2\gamma_2
\]
with $p_2$ in a standard proper rational $P_2<P$.
Rationality relative to $\Delta$ is rationality relative to $P\cap\Gamma$,
so $P_2$ is rational in $G$. The actual new factorization is
\[
                 g=(\epsilon\epsilon_2)p_2(\gamma_2\gamma).    \tag{D3}
\]
The order of both products is necessary in a noncommutative group.
Lemma~\ref{lem:closure} proves the small-coefficient assertions for the left
product, in the ambient coordinates. For the right product, the rational
lattice coordinates of $\gamma_2$ relative to $\Delta$ have common standard
denominators in the original rational structure. The same lemma supplies a
single discrete overgroup containing its values, the old overgroup, and
$\Gamma$. It also supplies the standard period modulo $\Gamma$.
Thus (D3) is a representation by the lower-dimensional subgroup $P_2$,
contrary to minimality.

No choice of $\Delta$ was imposed in that minimality argument. Consequently
the chosen $p$ is totally equidistributed for every standard finite-index
$\Delta\leq P\cap\Gamma$, with the same factors. This is the required
quantifier order. The argument only produces standard constants for this
particular internal orbit and fixed standard geometry. It does not assert
a prescribed effective bound uniform over all internal orbits.
\end{proof}

\section{A finite combinatorial proof of the linearization contract}
We prove a finite assertion with constants depending only on $c,C$ and then
transfer. Take $N\geq1$, $0<e\leq1$, and suppose the hypotheses of
Theorem~\ref{thm:F} hold in finite sets. Put $M=\lceil1/e\rceil$ and
$Q=100N+1$. Round each $\xi_i(h)$ in the circle to $y_i(h)\in M^{-1}\Z/\Z$,
with error at most $1/(2M)\leq e/2$. Since $Me\leq2$, the approximate sum
can take only a bounded number, depending on $C$, of values in this grid.
For at least $c_1N^3$ of the original quadruples its rounded value is a
fixed $a_0$ in that grid. Here $c_1>0$ depends only on $c,C$.

Embed $[-N,N]$ in $B=\Z/Q\Z$ and identify the target grid with $A=\Z/M\Z$.
In $B\times A$ form the four sets
\[
 A_1=\{(h,y_1(h))\},\quad A_2=\{(h,y_2(h))\},\quad
 A_3=\{(-h,y_3(h))\},\quad A_4=\{(-h,y_4(h))\}.             \tag{F1}
\]
The signed first coordinates are essential. The selected quadruples now
give $c_1N^3$ representations of $(0,a_0)$ in $A_1+A_2+A_3+A_4$.
The initial additive relation has no wraparound, since $4N<Q$.
Fourier inversion and Holder's inequality in this finite abelian group give
\[
 r_{A_1+A_2+A_3+A_4}(0,a_0)
          \leq\prod_{i=1}^4 E(A_i)^{1/4},\qquad
 E(A_i)=|\{x_1+x_2=x_3+x_4:x_a\in A_i\}|.
\]
For completeness, with unnormalized Fourier transforms the inverse formula
has factor $|B\times A|^{-1}$; Holder bounds the sum of four transforms
by the fourth moments, each of which equals $|B\times A|E(A_i)$.
Those factors cancel. Each $E(A_i)\leq |H|^3\leq27N^3$.
It follows that $E(A_1)\geq c_2N^3$ for a positive $c_2=c_2(c,C)$.
The same reasoning can isolate any $A_i$; reflecting its first coordinate
back if necessary gives the graph of the specified $y_i$.
We henceforth call that function $y$ and its graph $A_1$.

Apply [GT08, Theorems 5.2 and 5.3, p.95] to $A_1$ in the actual finite
group $B\times A$. Their hypotheses follow from the energy bound and
$|A_1|\leq3N$. They give a subset $\Gamma'$ of its graph and a constant
$K_0=K_0(c,C)$ such that
\[
            |\Gamma'|\geq N/K_0,
            \qquad |9\Gamma'-8\Gamma'|\leq K_0N.             \tag{F2}
\]
We use these published theorems themselves, not the special correlation
premise in [GT08, Proposition 5.4].

\subsection*{Graph refinement with the actual cyclic target}
Let $V=\{a\in A:(0,a)\in8\Gamma'-8\Gamma'\}$.
The map $\Gamma'\times V\to9\Gamma'-8\Gamma'$ given by addition is
injective: equality determines the first graph coordinate, then its unique
second coordinate, and finally $a$. Hence $|V|\leq K_0^2$.
There is a set $S_A$ of at most $O(1+\log K_0)$ characters of $A$ for which
\[
        V\cap\{a:\|\chi(a)\|_\T<1/4\ \forall\chi\in S_A\}=\{0\}.
\]
Indeed, for any nonzero $a$ a uniformly random character takes values
uniformly among the roots of unity of its order; at least one third of those
values have distance at least $1/4$ from zero. Choose enough independent
characters that the union bound over $V\setminus\{0\}$ is less than one.
This proves the required separation without identifying $A$ with the dual
of $B$, an identification that is unavailable here.

Partition each coordinate of the character vector on $A$ into intervals of
length at most $1/100$. A cell contains a subset $\Gamma''\subseteq\Gamma'$
of size $c_3N$. In any element of $8\Gamma''-8\Gamma''$, every selected
character of the second coordinate has a real representative of absolute
value less than $16/100<1/4$ when its first coordinate is zero.
It follows that the only element of this difference set with first
coordinate zero is $(0,0)$. Therefore $4\Gamma''-4\Gamma''$ is a graph.
Write $H''$ for the projection of $\Gamma''$, and let
\[
 \zeta:2H''-2H''\longrightarrow A
\]
be the function whose graph is $2\Gamma''-2\Gamma''$.
It is an exact Freiman homomorphism on this domain: two sums of two of its
graph points lie in $4\Gamma''-4\Gamma''$, where the first coordinate
determines the second. Also $\zeta(0)=0$. Thus
$\zeta(x+y)=\zeta(x)+\zeta(y)$ whenever $x,y,x+y$ lie in its domain.
This proves precisely the generalization of [GT08, Lemma 9.2] needed here.

\subsection*{A proper progression and an explicit linear bracket formula}
The set $H''$ has density at least $c_3/101$ in $B$.
By [GT08, Lemma 6.4, pp.97--98] there is a bounded set of characters $S$
with $B(S,1/4)\subseteq2H''-2H''$.
Add the faithful character $x\mapsto x/Q$ to $S$, and shrink the radius
to $\rho=1/16$. Write $d=|S|$, a bounded integer.
[GT08, Lemma 10.4, p.121] supplies a proper centered progression
\[
 P=\Bigl\{\sum_{a=1}^r l_av_a: |l_a|<L_a\Bigr\}\subseteq B(S,\rho),
              \qquad |P|\geq c_4 Q,
\]
with $r\leq d$, and linearly independent vectors
$\Phi(v_a)=(\{\chi(v_a)\})_{\chi\in S}$.
The coset subgroup in that lemma is zero because $S$ contains a faithful
character. Its proof also gives
\[
 \Phi\left(\sum_a l_av_a\right)=\sum_a l_a\Phi(v_a)\quad\hbox{on }P. \tag{F3}
\]
This is the lift of a lattice progression inside the real cube, not an
unjustified inference from reduction modulo one. Delete inactive generators
with no nonzero integer coefficient allowed; all remaining $v_a$ lie in $P$.
Every partial sum along a coordinate path to a point of this centered box
stays in $P$. The local additivity of $\zeta$ consequently gives
\[
          \zeta\left(\sum_a l_av_a\right)=\sum_a l_a\zeta(v_a).
\]
Choose real lifts $b_a$ of $\zeta(v_a)$ and real dual vectors $u_a$ to the
independent vectors $\Phi(v_a)$ on their span. Using (F3), define
$b_\chi=\sum_a b_a(u_a)_\chi$. Then, for every $p\in P$,
\[
                 \zeta(p)=\sum_{\chi\in S}b_\chi\{\chi(p)\}\pmod1. \tag{F4}
\]
This is a linear calculation. It uses neither the quadratic-extension
theorem nor the quadratic classification lemma in [GT08]. The author notice
concerning that extension theorem is therefore not a new assumption here.

\subsection*{Translation, fixed carries, and the Freiman restriction}
Averaging translates of $P$ in $B$ gives an $x\in B$ with
$|H''\cap(x+P)|\geq |H''||P|/Q\geq c_5N$.
For $h=x+p$ and $h_0=x+p_0$ in this intersection, the graph relation gives
$y(h)-y(h_0)=\zeta(p-p_0)$, since $\Gamma''-\Gamma''$ is contained in
$2\Gamma''-2\Gamma''$. Both $p,p_0$ lie in the domain of $\zeta$.
Its Freiman identity applied to $(p-p_0)+p_0=p+0$ gives
\[
                 y(h)=\delta_x+\zeta(p).
\]
For each $\chi\in S$, write $\chi(h)=\alpha_\chi h$ on the cyclic group.
The difference
\[
 \{\chi(h-x)\}-\{\chi(h)\}+\{\chi(x)\}
\]
is an integer in a fixed set of at most three values. Restrict to a cell on
which all these carries are fixed. This costs a factor at most $3^d$ and,
by (F4), gives
\[
                y(h)=\delta+\sum_{\chi\in S}b_\chi\{\alpha_\chi h\}.
\]
Finally partition each $\{\alpha_\chi h\}$ into a fixed number of intervals
of length at most $1/10$, and retain a largest cell $H'$.
For any integer relation $h_1+h_2=h_3+h_4$ in $H'$, the corresponding sum
of four fractional parts, with signs $+,+,-,-$, is an integer of absolute
value at most $1/5<1$. It is therefore zero in $\R$ for each character.
The real weighted sum is an exact Freiman homomorphism on $H'$.
Undoing the rounding changes the displayed value by at most $e/2$.
All density losses depend only on $c,C$ and all sets and choices are finite.
This proves the finite assertion, and transfer with $e=N^{-j}$ proves
Theorem~\ref{thm:F}. A bounded rank can be made a fixed standard rank by
passing to its ultrapower value or by adding zero terms.

\section{Consumer boundary and related-use scope}
Theorem~\ref{thm:D} supplies exactly the one-variable factorization consumed
by the Sections 9--11 repair. That repair still must derive its nonzero coset
integral and its degree-rank horizontal inclusion for the actual Taylor
quotients. In particular, it may multiply a character by a standard
denominator to kill the rational periodic contribution; it may not silently
erase that contribution. No such application is accepted by this packet.

Theorem~\ref{thm:F} applies separately at each of the finitely many required
indices $j$, including $j=1$, and to any chosen one of the four functions.
The number of applications is standard. Dense refinements and all losses
remain standard only when that finite iteration is recorded by the consumer.
The theorem supplies no global Freiman extension beyond $H'$.

The related-use audit is bounded by uses of the single-scale equivalence,
properness in D.4, composition of factorization factors, deletion of rational
Taylor contributions, and the graph/Bohr/bracket joins in F.1. The
multi-variable D.5 and its consumers in Appendix E require their own exact
application review. They are not inferred from the one-variable theorem here.
Likewise, this packet does not accept Gaussian's pending Sections 9--11
argument or any downstream assembled inverse theorem.
\begin{thebibliography}{9}
\bibitem{GTZ}[GTZ] B. Green, T. Tao and T. Ziegler,
An inverse theorem for the Gowers $U^{s+1}[N]$-norm,
\emph{Annals of Mathematics} 176 (2012), 1231--1372.
DOI: 10.4007/annals.2012.176.2.11.
\bibitem{GT12}[GT12] B. Green and T. Tao,
The quantitative behaviour of polynomial orbits on nilmanifolds,
\emph{Annals of Mathematics} 175 (2012), 465--540.
DOI: 10.4007/annals.2012.175.2.2.
\bibitem{GT14}[GT14] B. Green and T. Tao,
On the quantitative distribution of polynomial nilsequences---erratum,
\emph{Annals of Mathematics} 179 (2014), 1175--1183.
DOI: 10.4007/annals.2014.179.3.8.
\bibitem{GT08}[GT08] B. Green and T. Tao,
An inverse theorem for the Gowers $U^3(G)$ norm,
\emph{Proceedings of the Edinburgh Mathematical Society} 51 (2008), 73--153.
DOI: 10.1017/S0013091505000325.
\end{thebibliography}
\end{document}
